\PropositionE{5.14}
{If a proposition follows from another, then the
latter says more than the former, the former less
than the latter.}
% -----File: 111.png---


\PropositionE{5.141}
{If $p$ follows from $q$ and $q$ from $p$ then they are
one and the same proposition.}


\PropositionE{5.142}
{A tautology follows from all propositions: it
says nothing.}


\PropositionE{5.143}
{Contradiction is something shared by propositions,
which \emph{no} proposition has in common with
another. Tautology is that which is shared by
all propositions, which have nothing in common
with one another.

Contradiction vanishes so to speak outside,
tautology inside all propositions.

Contradiction is the external limit of the propositions,
tautology their substanceless centre.}


\PropositionE{5.15}
{If $T_{r}$ is the number of the truth-grounds of the
proposition ``$r$'', $T_{rs}$ the number of those truth-grounds
of the proposition ``$s$'' which are at the
same time truth-grounds of ``$r$'', then we call the
ratio $T_{rs} : T_{r}$ the measure of the \emph{probability} which
the proposition ``$r$'' gives to the proposition ``$s$''.}


\PropositionE{5.151}
{Suppose in a schema like that above in No.~\PropERef{5.101}
$T_{r}$ is the number of the ``T'''s in the proposition
$r$, $T_{rs}$ the number of those ``T'''s in
the proposition $s$, which stand in the same columns
as ``T'''s of the proposition $r$; then the proposition
$r$ gives to the proposition $s$ the probability
$T_{rs} : T_{r}$.}


\PropositionE{5.1511}
{There is no special object peculiar to probability
propositions.}


\PropositionE{5.152}
{Propositions which have no truth-arguments
in common with one another we call independent.
\enlargethispage{-3pt} % force next paragraph to the next page

Independent propositions (\exempliGratia\ any two elementary
propositions) give to one another the probability~$\frac{1}{2}$.

If $p$ follows from $q$, the proposition $q$ gives
to the proposition $p$ the probability~1. The
certainty of logical conclusion is a limiting case
of probability.

(Application to tautology and contradiction.)}


\PropositionE{5.153}
{A proposition is in itself neither probable nor
% -----File: 113.png---
improbable. An event occurs or does not occur,
there is no middle course.}


\PropositionE{5.154}
{In an urn there are equal numbers of white
and black balls (and no others). I draw one
ball after another and put them back in the
urn. Then I can determine by the experiment
that the numbers of the black and white balls
which are drawn approximate as the drawing
continues.

So \emph{this} is not a mathematical fact.

If then, I say, It is equally probable that
I should draw a white and a black ball, this
means, All the circumstances known to me (including
the natural laws hypothetically assumed)
give to the occurrence of the one event no more
probability than to the occurrence of the other.
That is they give---as can easily be understood
from the above explanations---to each the
probability~$\frac{1}{2}$.

What I can verify by the experiment is that
the occurrence of the two events is independent
of the circumstances with which I have no closer
acquaintance.}


\PropositionE{5.155}
{The unit of the probability proposition is: The
circumstances---with which I am not further acquainted---give
to the occurrence of a definite event
such and such a degree of probability.}


\PropositionE{5.156}
{Probability is a generalization.

It involves a general description of a propositional
form. Only in default of certainty do we
need probability.

If we are not completely acquainted with a fact,
but know \emph{something} about its form.

(A proposition can, indeed, be an incomplete
picture of a certain state of affairs, but it is always
\emph{a} complete picture.)
% -----File: 115.png---

The probability proposition is, as it were, an
extract from other propositions.}
